% Direction artistique : édition scientifique, octobre 2026.
% Georgia pour les titres et les chiffres ; Calibri pour la lecture et les annotations.
\setsansfont{Calibri}[Path=../../../docs/collaboration/local/da_fonts/,UprightFont=calibri.ttf,
  BoldFont=calibrib.ttf,ItalicFont=calibrii.ttf,BoldItalicFont=calibriz.ttf]
\newfontfamily\editorial{Georgia}[Path=../../../docs/collaboration/local/da_fonts/,UprightFont=georgia.ttf,
  BoldFont=georgiab.ttf,ItalicFont=georgiai.ttf,BoldItalicFont=georgiaz.ttf]
\definecolor{ardoise}{HTML}{252C32}
\definecolor{acier}{HTML}{68747C}
\definecolor{encre}{HTML}{252C32}
\definecolor{gris}{HTML}{626A70}
\definecolor{accent}{HTML}{A52D39}
\definecolor{clair}{HTML}{FAFAF9}
\definecolor{filet}{HTML}{DADBD9}
\definecolor{bleuclair}{HTML}{F1F2F1}
\definecolor{nexialog}{HTML}{223E55}
\setbeamercolor{normal text}{fg=encre,bg=white}
% Navigation Beamer : sections et une pastille par frame (pas par overlay).
\setbeamercolor{section in head/foot}{fg=white,bg=nexialog}
\setbeamercolor{mini frame}{fg=white,bg=nexialog}
\setbeamerfont{section in head/foot}{size=\fontsize{5.8}{7}\selectfont,series=\mdseries}
\setbeamertemplate{section in head/foot shaded}[default][45]
\setbeamertemplate{mini frames}[default]
\setbeamertemplate{headline}{%
  \begin{beamercolorbox}[wd=\paperwidth]{section in head/foot}%
    \vspace{.7mm}%
    \insertnavigation{\paperwidth}%
    \vspace{.7mm}%
  \end{beamercolorbox}%
  \vspace{.5mm}%
  \hspace*{7mm}{\color{filet}\rule{\dimexpr\paperwidth-14mm}{.35pt}}\par
}
\setbeamerfont{framesubtitle}{size=\small,series=\mdseries}
\setbeamercolor{framesubtitle}{fg=gris}
\setbeamertemplate{blocks}[default]
\setbeamerfont{block title}{family=\editorial,size=\normalsize,series=\mdseries}
\setbeamerfont{block title alerted}{family=\editorial,size=\normalsize,series=\mdseries}
\setbeamerfont{block title example}{family=\editorial,size=\normalsize,series=\mdseries}
\setbeamercolor{block title}{fg=encre,bg=white}
\setbeamercolor{block body}{fg=encre,bg=white}
\setbeamercolor{block title alerted}{fg=accent,bg=white}
\setbeamercolor{block body alerted}{fg=encre,bg=white}
\setbeamercolor{block title example}{fg=encre,bg=white}
\setbeamercolor{block body example}{fg=encre,bg=white}
\setbeamertemplate{itemize item}{\raisebox{.25ex}{\color{accent}\rule{3pt}{.5pt}}}
\setbeamertemplate{itemize subitem}{\raisebox{.25ex}{\color{gris}\rule{2pt}{.4pt}}}
% Miniature native du schéma directeur ; numéros identiques aux étapes du schéma complet.
\def\methodeactive{}
\def\methodetitre{}
\newcommand{\suivimethode}[2]{\gdef\methodeactive{#1}\gdef\methodetitre{#2}}
\newcommand{\minimethode}{%
  \begin{tikzpicture}[x=1mm,y=1mm,baseline=(current bounding box.north),
    rep/.style={draw=nexialog!35,fill=white,text=gris,inner sep=0pt,align=center,
      font=\fontsize{4.6}{5.2}\selectfont},
    actif/.style={draw=nexialog,fill=nexialog,text=white,font=\fontsize{4.6}{5.2}\selectfont\bfseries},
    lien/.style={-{Stealth[length=.7mm,width=.6mm]},draw=nexialog!45,line width=.3pt}]
    % disposition du schéma directeur : 1 | 2 | 3-4-5 | 6-7 | 8
    \foreach \n/\x/\y/\w/\h/\lab in {%
        1/3/-5/6/10/{\textbf{1}\\État},
        2/11/-5/8/10/{\textbf{2}\\Canaux},
        3/22.5/-1.6/11/2.8/{\textbf{3} Fréquence},
        4/22.5/-5/11/2.8/{\textbf{4} Sévérité},
        5/22.5/-8.4/11/2.8/{\textbf{5} Cascade},
        6/34/-2.1/9/4/{\textbf{6} Perte},
        7/34/-7.1/9/4/{\textbf{7} Capital},
        8/34/-11.8/9/3/{\textbf{8} Entité}}{%
      \def\estactif{0}%
      \foreach \a in \methodeactive {\ifnum\n=\a\relax\xdef\estactif{1}\fi}%
      \ifnum\estactif=1
        \node[rep,actif,minimum width=\w mm,minimum height=\h mm] (m\n) at (\x,\y) {\lab};
      \else
        \node[rep,minimum width=\w mm,minimum height=\h mm] (m\n) at (\x,\y) {\lab};
      \fi}
    \draw[lien] (m1)--(m2);
    \foreach \n in {3,4,5}{\draw[lien] (m2.east|-m\n)--(m\n.west);}
    \draw[lien] (m3.east)--([yshift=.6mm]m6.west);
    \draw[lien] (m4.east)--([yshift=.9mm]m7.west);
    \draw[lien] (m5.east)--([yshift=-.9mm]m7.west);
    \draw[lien] (m6)--(m7);
    \draw[lien] (m7)--(m8);
    \node[anchor=south east,text=nexialog,font=\fontsize{5.2}{6}\selectfont\itshape,inner sep=0pt]
      at (38.5,.9) {Démarche · \methodetitre};
  \end{tikzpicture}%
}
\makeatletter
\setbeamertemplate{frametitle}{%
  \vspace*{1mm}%
  \ifx\methodeactive\@empty\def\margeright{7mm}\else\def\margeright{49mm}\fi
  \begin{beamercolorbox}[wd=\paperwidth,leftskip=7mm,rightskip=\margeright]{frametitle}%
    \ifx\methodeactive\@empty\else
      \begin{tikzpicture}[remember picture,overlay]
        \node[anchor=north east,inner sep=0pt] at ([xshift=-7mm,yshift=-6.2mm]current page.north east) {\minimethode};
      \end{tikzpicture}%
    \fi
    \ifx\methodeactive\@empty\else\vbox to 15mm\bgroup
      \hsize=\dimexpr\paperwidth-56mm\relax\leftskip=0pt\rightskip=0pt\fi
    {\editorial\mdseries\fontsize{14}{16}\selectfont\color{encre}\insertframetitle\par}%
    \ifx\insertframesubtitle\@empty\else
    \vspace{1mm}{\small\mdseries\color{gris}\insertframesubtitle\par}\fi
    \ifx\methodeactive\@empty\else\vfil\egroup\fi
  \end{beamercolorbox}%
  \vspace{1mm}%
}
\makeatother
\setbeamertemplate{footline}{%
  \hspace*{7mm}{\color{filet}\rule{\dimexpr\paperwidth-14mm}{.35pt}}\par
  \vspace{.6mm}\hspace*{7mm}%
  \makebox[\dimexpr\paperwidth-14mm][l]{%
    {\fontsize{5.5}{7}\selectfont\color{gris}KÉLIAN KADDOURI\quad /\quad ENSAE · NEXIALOG}%
    \hfill{\editorial\fontsize{8}{9}\selectfont\color{accent}\insertframenumber}}%
  \par\vspace{1.4mm}}
% Statut d'une diapositive : ce que le travail ajoute (rouge) ou reprend (gris).
\newcommand{\badge}[2]{\tikz[baseline=(b.base)]\node[fill=#1,text=white,rounded corners=.6mm,inner sep=.9mm,
  font=\fontsize{5.6}{6.5}\selectfont\bfseries] (b) {#2};}
\newcommand{\ajout}{\badge{accent}{AJOUT}\enspace}
\newcommand{\etabli}{\badge{acier}{ÉTABLI}\enspace}
% Question de transition vers la diapositive suivante.
\newcommand{\question}[1]{\par\vspace{2mm}{\small\color{accent}\itshape$\rightarrow$\enspace #1\par}}
\renewcommand{\chiffre}[2]{{\editorial\fontsize{17}{19}\selectfont\color{accent}#1}\par\vspace{1mm}{\scriptsize #2\par}}
\newcommand{\repere}[2]{{\color{accent}\bfseries #1}\enspace\textbf{#2}}
\tikzset{etape/.append style={rounded corners=0pt,fill=white},pilier/.append style={fill=white}}
